\section{Design alternatives that were considered}

Several decisions in the design have an obvious alternative. Table~\ref{tab:alt} records what was chosen and why. The reasoning matters more than the choice, because it shows what the constraints were.

\begin{table}[ht]
\centering
\small
\begin{tabularx}{\textwidth}{p{2.9cm}p{3.4cm}X}
\toprule
Question & Chosen & Reason, and what was rejected \\
\midrule
Integration point & kernel crypto API providers & Every program can use the board unchanged. A private character device with its own ioctl interface would have needed a client library and would have been invisible to the kernel's own users. \\
Board side & user space daemon over FunctionFS & Ordinary C code, cross compiled and debugged with ordinary tools. A second kernel module on the board would have needed a kernel rebuild and re-flash, and the board had no recovery path. \\
Message framing & one request per bulk transfer & Boundaries are given by USB itself, so no length prefix scanning is needed. A byte stream over a serial style endpoint would have needed its own framing and resynchronisation. \\
Flow control & window of tagged commands & Overlaps transfer and compute. Stop and wait left both sides idle half of the time (Chapter~\ref{ch:test} measures this). \\
Loss handling & host timer plus board reply cache & Retransmission is safe because the board never executes a tag twice. Relying on USB's own retry would not cover a daemon that dropped a request it had read. \\
Ciphers & asynchronous skcipher & Matches how hardware drivers are written and lets many requests share the engine. A synchronous interface would have blocked the caller for the full round trip. \\
Hashes & synchronous shash and asynchronous ahash & The shash interface is what most in-kernel users call. The ahash interface shows the asynchronous style for hashing. Both use one engine. \\
GCM on board & assembled from CTR, ECB and software GHASH & The board kernel has no AEAD socket support. Rebuilding it was rejected as too risky. \\
Key storage & key sent with every command & The board holds no key state, so a reboot or retransmission needs no recovery logic. The cost is that keys cross the cable (Chapter~\ref{ch:security}). \\
\bottomrule
\end{tabularx}
\caption{Design decisions and the alternatives rejected.}\label{tab:alt}
\end{table}

\section{Failure model}

A design is only as good as the failures it expects. The project assumes the following can happen, and each has a mechanism that deals with it.

\begin{table}[ht]
\centering
\begin{tabularx}{\textwidth}{lXX}
\toprule
Failure & Detected by & Handled by \\
\midrule
Request packet lost on the board & no reply within one second & host resends the same tag, up to four sends \\
Reply packet lost & same timeout & the board answers the repeat from its reply cache \\
Duplicate reply arrives & no command with that tag is in flight & ignored \\
Board answers with an error & status field in the reply & error number returned to the caller \\
Garbled or unknown message & magic and length checks on both sides & dropped (the host will retry) \\
Cable pulled & URB completes with a shutdown status; disconnect is called & everything pending is completed with \code{-ENODEV} \\
No board attached & the board list is empty & new requests return \code{-ENODEV} \\
Module reloaded while the board stays up & random starting tag & stale cached replies are not matched \\
Memory allocation fails & \code{NULL} return & \code{-ENOMEM} completes the request \\
Board daemon stops answering & no reply after all retries & \code{-ETIMEDOUT} \\
\bottomrule
\end{tabularx}
\caption{Failures and their handling.}
\end{table}

\section{Who owns what}

The responsibilities of the three components are deliberately narrow.

\begin{itemize}
  \item \textbf{The module} owns the interface to the kernel (algorithms, requests, error codes), the cutting of requests into commands, windowing, matching and retransmission.
  \item \textbf{The protocol header} owns the format of every message and the limits. Neither side defines a number that the other must guess.
  \item \textbf{The daemon} owns execution: it maps an operation code to an AF\_ALG call, keeps hash sessions, remembers recent replies, and hides how the board does the work.
\end{itemize}

Because the module treats the board as a black box that executes tagged commands, the same module would work unchanged with a daemon that used a different engine, as long as it honours the protocol.

\section{Concurrency at a glance}

Figure~\ref{fig:conc} shows the threads and execution contexts in play on each side.

\begin{figure}[ht]
\centering
\begin{tikzpicture}[node distance=5mm and 8mm]
  \node[font=\small\bfseries] at (2.5,4.4) {PC (module)};
  \node[box,fill=blue!8,minimum width=4.6cm] (callers) at (2.5,3.4) {many callers\\(processes, kernel users)};
  \node[box,fill=orange!15,minimum width=4.6cm] (wq) at (2.5,1.9) {one ordered workqueue\\pump, done, fail, timeout};
  \node[box,fill=gray!15,minimum width=4.6cm] (irq) at (2.5,0.4) {URB completion (softirq)\\receive and send complete};
  \draw[arr] (callers) -- (wq);
  \draw[arr] (wq) -- (irq);
  \draw[arr] (irq.east) -- ++(0.5,0) |- (wq.east);

  \node[font=\small\bfseries] at (9.5,4.4) {Board (daemon)};
  \node[box,fill=blue!8,minimum width=4.6cm] (m) at (9.5,3.4) {main thread\\ep0 events};
  \node[box,fill=green!18,minimum width=4.6cm] (w) at (9.5,1.9) {worker: serving loop};
  \node[box,fill=orange!15,minimum width=2.2cm] (h) at (8.2,0.4) {harvester};
  \node[box,fill=orange!15,minimum width=2.2cm] (wt) at (10.8,0.4) {writer};
  \draw[arr] (m) -- (w);
  \draw[arr] (h) -- (w);
  \draw[arr] (w) -- (wt);
\end{tikzpicture}
\caption{Execution contexts on each side.}\label{fig:conc}
\end{figure}

On the PC, the ordered workqueue means that only one work item runs at a time per module. This removes whole classes of races between, for example, the code that emits new chunks and the code that completes a request. Callers and the URB completion handlers still run concurrently with the workqueue, and the spinlock \code{d->lock} covers them.
